\documentclass[times,final,authoryear]{elsarticle}

\usepackage{medima}
\usepackage{framed,multirow}
\usepackage{graphicx}
\usepackage{booktabs}
\usepackage{subcaption}
\usepackage{tabularx}
\usepackage{threeparttable}

\usepackage{amssymb}
\usepackage{amsmath}
\usepackage{latexsym}
\usepackage[left]{lineno}

\usepackage{tikz}
\usetikzlibrary{
    calc,
    positioning,
    arrows.meta,
    shapes.geometric,
    shapes.symbols
}
\usepackage{url}
\usepackage{xcolor}

\usepackage{hyperref}
\hypersetup{hidelinks}

\definecolor{newcolor}{rgb}{.8,.349,.1}

\journal{Agricultural Systems}

\makeatletter
\gdef\@jnlurl{https://www.sciencedirect.com/journal/agricultural-systems}
\def\@issn{Agricultural Systems}

\def\elsjnllogoboxb#1#2{\minipage[b][#2][b]{#1}\mbox{}\endminipage}
\makeatother

\begin{document}

\verso{Martins \textit{et~al.}}

\begin{frontmatter}

\title{Assessing Context-Conditioned Soybean Sowing Date and Supplemental Irrigation Decisions Using DSSAT and Reinforcement Learning}

\author[2,6,7,8]{Marcella  Scoczynski Ribeiro \snm{Martins}} %0000-0002-5716-4968
\author[1]{Gabrielly de Queiroz \snm{Pereira}} %orcid: https://orcid.org/0000-0003-2663-0783 gqpereira@uepg.br
\author[2]{Hilkija Gaïus \snm{Tosso}} %orcid https://orcid.org/0000-0003-1109-6813 hgaius@hotmail.com
\author[1,8]{Maria Salete Marcon Gomes \snm{Vaz}} %orcid https://orcid.org/0000-0001-9172-1863 salete@uepg.br
\author[1,8]{José Carlos Ferreira da \snm{Rocha}} %orcid https://orcid.org/0000-0002-4050-281X jrocha@uepg.br
\author[4,5,6]{Viviane M. \snm{Gomes Pacheco}} %0000-0003-3693-4377 viviane.gomes@ifg.edu.br
\author[4]{Cloves Gonçalves \snm{Rodrigues}} %0000-0003-0140-9847 cloves@pucgoias.edu.br
\author[3]{Antonio~Paulo~\snm{Coimbra}} %0000-0003-3780-2405 acoimbra@isr.uc.pt
\author[3,5,6]{Wesley Pacheco \snm{Calixto}\corref{cor1}} %0000-0002-1928-4432 wesley.pacheco@isr.uc.pt
\cortext[cor1]{Corresponding author.}
\ead{wesley.pacheco@isr.uc.pt}

\address[1]{State University of Ponta Grossa, Ponta Grossa, Parana/Brazil}
\address[2]{Federal University of Technology - Parana, Brazil}
\address[3]{Institute of Systems and Robotics, Coimbra University, Coimbra/Portugal}
\address[4]{Polytechnic and Arts School, Pontifical Catholic University of Goias, Goiania, Goias/Brazil PC~74.605-010}
\address[5]{Electrical, Mechanical \& Computer Engineering School, Federal University of Goias, Goiania, Goias/Brazil PC~74.605-010}
\address[6]{Experimental \& Technological Research and Study Group, Federal Institute of Goias, Goiania, Goias,  Brazil}
\address[7]{Federal University of Technology - Parana, Graduate Program in Electrical Engineering, Ponta Grossa, Brazil}
\address[8]{State University of Ponta Grossa, Postgraduate Program in Applied Computing, Ponta Grossa, Parana/Brazil}



\begin{abstract}
\textbf{CONTEXT:} Sowing date and supplemental irrigation are coupled decisions in predominantly rainfed soybean systems and are chosen before seasonal weather is known. A context-conditioned rule must be evaluated against comparators that receive the same information and simulation budget, on seasons not used to build it.

\textbf{OBJECTIVE:} To formulate soybean sowing date and irrigation-rule selection as a one-step contextual decision problem coupled to DSSAT-CROPGRO-Soybean, and to test whether a policy trained with Proximal Policy Optimization (PPO) on information available on September 1 adds value over an optimized fixed rule and a simple contextual rule in Castro, Paraná, Brazil.

\textbf{METHODS:} PPO selected sowing date and three parameters of a causal irrigation rule from six pre-season variables (station rainfall, temperature and radiation before the decision date, and the Oceanic Ni\~no Index). Each action was evaluated by a full DSSAT run with raw simulated yield and a soil-water spin-up. Seasons were split chronologically into training (2006--2016), validation (2017--2020) and a frozen test block (2021--2024). Five training seeds were compared, season by season, with an optimized fixed rule, a nearest-neighbour contextual rule, constructed schedules and a retrospective oracle evaluated on the same grid of 602 candidate rules.

\textbf{RESULTS:} PPO exceeded the optimized fixed rule on the training seasons (objective 1.156 vs 1.130) and on validation (1.226 vs 1.175). On the test block the order was reversed: mean yield was 4,061 kg ha$^{-1}$ for PPO and 4,904 kg ha$^{-1}$ for the optimized fixed rule, a paired difference of $-843$ kg ha$^{-1}$, negative in all four seasons. Two mechanisms explained the loss: station temperature records were missing for the whole pre-decision window in 2021 and 2022, so the imputed context fell 4--7 standard deviations outside the training range and the policies delayed sowing to late November; and the four validation seasons required no irrigation, so checkpoint selection retained policies that did not irrigate in the 2024 drought, when the fixed rule applied 75 mm and produced 2,564 kg ha$^{-1}$ more than rainfed management.

\textbf{CONCLUSIONS:} With 11 training seasons, the context-conditioned policy learned season-specific actions that did not transfer to new seasons, while a fixed rule with an in-season irrigation trigger, selected on all development seasons, captured 67\% of the oracle gain over the rainfed calendar. Data quality of pre-season inputs and the hydrological composition of the validation block determined out-of-sample performance more than the learning algorithm.

\end{abstract}

\begin{keyword}
Soybean production \sep Reinforcement learning\sep DSSAT Simulations \sep Proximal Policy Optimization \sep Decision analysis \sep Water allocation
\end{keyword}

\end{frontmatter}

\linenumbers

%===============================================================
\section{Introduction}
\label{introduction}

In soybean production in southern Brazil, water deficit is the main cause of the yield gap and irrigation is among the remedies proposed to close it \citep{sentelhas2015yieldgap}. A farmer chooses, once per season, a sowing date inside a window of roughly three months and, where equipment exists, a rule for supplemental irrigation. The outcome of that choice depends on the interaction between seasonal weather, soil water-holding capacity, cultivar phenology, and the amount of water available for irrigation, none of which is fully known at sowing. A fixed sowing calendar ignores this interaction. A decision rule that conditions sowing date and water use on information available before the season could, in principle, exploit it. Whether such a rule delivers a measurable benefit over a well-chosen calendar, and under which water restrictions within the Castro production context, is a question of system behaviour that agricultural systems modelling was developed to address \citep{jones2017history}. It is also a question in which the value of pre-season climate information is itself part of the analysis \citep{hansen2002climate}.

Agricultural management requires selecting actions whose consequences depend on interacting environmental and biological conditions. Sowing date, irrigation, cultivar traits, soil properties, and weather jointly affect crop development and final yield; for this reason, management design is better treated as a decision and optimization problem than as yield prediction under a fixed strategy \citep{tao2022optimizing, Xiao2024}. Recent work in agricultural systems has also emphasized that digital decision systems are not defined only by predictive accuracy: the objective function, its weights, constraints, and trade-off rules shape what the system is designed to optimize \citep{carolan2026objective}. This point is directly relevant to irrigation and sowing decisions, because increasing yield, reducing water use, limiting risk, and maintaining operational feasibility can lead to different preferred actions. In soybean production, sowing date is particularly important because it determines the environmental conditions experienced during critical phenological stages and can also affect the feasibility and productivity of subsequent crops within the same agricultural year \citep{debangshi2025plantingdates, bigolin2025latesoybean, zhang2021modissoybean}.

Process-based crop simulation models provide a suitable computational basis for this type of analysis. Unlike purely empirical approaches that estimate yield directly from historical observations, these models represent physiological and environmental processes governing crop growth and development. The Decision Support System for Agrotechnology Transfer (DSSAT) is one of the most widely used crop-modeling platforms and enables the simulation of crop responses to weather, soil, genotype, and management conditions \citep{ovando2018dssatsoybean, singh2024dssatsoybean}. Previous soybean applications have examined responses to temperature, precipitation, sowing date, and weather-data sources \citep{battisti2017soybeanmodels, battisti2018climatesensitivity, ovando2018dssatsoybean}, typically for a predefined set of treatments. As the number of decision variables grows, exhaustive evaluation becomes less efficient. Applications of an established crop model to a decision problem also presuppose that the model has been calibrated and tested for the case in question. When soil and cultivar parameters are generic, the simulated response carries an unquantified systematic error that any downstream optimization inherits.

Optimization has become a common framework for agricultural water-resource management because it allows management scenarios to be compared under constraints and competing objectives. Recent reviews of agricultural water and reservoir management describe the use of simulation--optimization models, multiobjective optimization, and machine-learning surrogate models to reduce computational cost when evaluating management alternatives \citep{kumari2026optimization}. In parallel, Agriculture 5.0 and drone-as-a-service systems are expanding the availability of data streams for crop monitoring, soil assessment, and automated farm operations \citep{adel2026drones}. These data streams can improve the context available to decision models, but they do not by themselves answer whether a context-aware policy outperforms a well-selected fixed rule under the same information and simulation budget.

Machine learning has been used to forecast soybean yield \citep{vonbloh2023mlsoybean, schwalbert2020soybeanforecast}. Forecasting, however, estimates the outcome expected under a given set of conditions or management assumptions, while management optimization requires choosing among alternative actions. Reinforcement learning addresses this second problem by learning a policy through repeated interaction with an environment, which may be a process-based crop simulator. In irrigation scheduling, \citet{kelly2024assessing} showed with AquaCrop-OSPy that deep RL provided its clearest advantage under strong water limitation, while optimized heuristic policies were difficult to outperform under wetter or more variable rainfall conditions. Using gym-DSSAT, \citet{balderas2025comparative} evaluated reinforcement learning algorithms for crop management tasks involving fertilization and irrigation. For planting-date adaptation, \citet{noa2026hybrid} combined a process-based crop model, a machine-learning emulator, and reinforcement learning to explore adaptive corn planting policies. The present study differs from these works by focusing on a joint pre-season sowing-date and irrigation-rule decision for soybean in Castro, with DSSAT executed directly for the evaluated management scenarios.

When the whole management strategy is fixed before the season, the decision problem has a single step and is a contextual optimization problem rather than a sequential control problem. Policy-gradient methods such as PPO remain applicable because they learn a stochastic mapping from seasonal context to a continuous action vector. Their use, however, must be judged against an optimized fixed rule, against context-free global optimizers, and against a simple contextual rule that receives the same information under the same simulation budget. In this study, PPO is one mechanism of context-conditioned policy search. The object of the study is the value of the context and of the management objective, not the algorithm alone.

Among the closest studies identified \citep{kelly2024assessing, balderas2025comparative, noa2026hybrid}, and in light of recent reviews and conceptual work in agricultural systems \citep{carolan2026objective, kumari2026optimization, adel2026drones}, we found no evaluation of a joint pre-season sowing-and-irrigation decision for soybean in which the comparators receive the same information and the same simulation budget as the learned policy, in which a simple contextual rule is included, and in which performance inside the training seasons is reported next to performance on later seasons. The novelty claimed here is restricted to these verifiable axes within the Castro case study. The central question is: \emph{under which conditions of climate, soil water storage, and water restriction does a contextual sowing-and-irrigation policy add value relative to a well-optimized fixed rule in Castro?} The question admits a null or negative answer, and such an answer is informative for the design of decision systems when its causes can be identified.

Three directional expectations were defined before the final analysis. First, the value of adaptation may interact with water availability and with the restriction on irrigation water, without a presumed monotonic direction. Second, the storage capacity of the soil may modulate the value of adaptation in a season-dependent way. Third, the value of pre-season information is related to its observed predictive skill for the coming season, measured on the development seasons before the policy is evaluated, without imposing an a priori ordering among climatology, pre-season observations, and seasonal climate indices. Each expectation has the unit of comparison of Section~\ref{subbaselines} and the pre-specified criterion of Section~\ref{subdesign}; the Conclusion reports the estimated direction and magnitude rather than confirmed or refuted hypotheses.

The study addresses four analytical objectives: i) to quantify the performance of a context-conditioned sowing-and-irrigation policy relative to an optimized fixed rule in Castro seasons not used for model development; ii) to determine whether the value of contextual adaptation persists when the policy is compared with a simple contextual rule receiving the same pre-season information and with constructed reference schedules; iii) to quantify how contextual adaptation modifies the productivity--water relation and the yield of the least favourable seasons; and iv) to identify the conditions under which pre-season context changes the selected management action, including the predictive skill of each context variable and the coverage of the training seasons.

The intended contribution is to quantify, for Castro, whether context-conditioned soybean sowing and irrigation-rule selection provides value beyond a well-selected fixed rule, and to identify why the answer differs between the seasons used for training and later seasons. The contribution is methodological and diagnostic. It documents a reproducible DSSAT--policy coupling with raw simulated yield, a comparator structure in which every rule is evaluated on the same simulated seasons, learning curves and checkpoint selection for five training seeds, and the effect of data quality in the pre-season context on the learned policy. Spatial transfer to other municipalities is outside the scope of this study.

% =====================================================================
\section{Materials and methods}
\label{metodologia}

%%============================================================
\subsection{System boundary and study environment}
\label{subsystem}

The system studied is a predominantly rainfed soybean field with optional supplemental irrigation, represented for Castro by one weather station, one soil profile, and one cultivar. Its spatial boundary is the field (one soil, one management); its temporal boundary is the growing season from sowing to physiological maturity, preceded by a soil-water spin-up that starts on August 1 of the sowing year (Section~\ref{subcropmodel}). The components and their interactions are: (i) the seasonal weather (rainfall, temperature, radiation), uncertain at sowing; (ii) the soil water balance, which converts rainfall and irrigation into water available to the crop and depends on the profile and on the antecedent moisture; (iii) the crop, whose phenology and yield formation respond to the timing of water and thermal stress relative to the sowing date; (iv) the management decision, taken once before the season, consisting of a sowing date within a window and a rule for supplemental irrigation; and (v) the water resource, bounded by a seasonal cap. The unit of decision is one season of one field. Nitrogen is assumed non-limiting (Section~\ref{subcropmodel}), so the conclusions do not extend to nutrient-constrained systems.

\emph{Decision instant and information horizons}: the operational decision date is fixed at $t_0=$ September 1 of the sowing year. The policy receives the context and emits the action before the sowing window opens on September 15. The admissible station information ends on August 31, and the admissible climate-index information is the last value published before $t_0$ (Section~\ref{subdecision}). The decision is not revised after $t_0$; the irrigation rule then uses daily weather during the season only through its prescribed causal schedule, and DSSAT simulates from the spin-up date to physiological maturity. Information reaches the decision maker at two moments: before $t_0$ (admissible in the context) and during the season (used by the simulator, by the causal irrigation rule, and, in a labelled retrospective diagnostic, by the oracle).

The environment is the Castro station (Paraná, Brazil, latitude $-24.78$, longitude $-50.05$, altitude 1,016 m), with a synthetic clayey profile, the generic maturity-group-5 cultivar \texttt{990005} of \texttt{SBGRO048.CUL}, plant population 30 m$^{-2}$, row spacing 45 cm, and sowing depth 5 cm. Table~\ref{tab_environment_a_current} documents the environment and, by showing zero field-validation observations, also states why the crop-model outputs are treated as a simulation experiment and not as a field prediction. The complete layer profile is exported by the code as \texttt{synthetic\_soil\_profile\_s2.csv}. No claim is made for spatial generalization beyond Castro.

\begin{table}[!h]
\centering
\scriptsize
\caption{Executable environment. Field validation metrics are not estimated because no field phenology, yield, or soil-water observations are available for this environment.}
\label{tab_environment_a_current}
\begin{tabular}{p{2.0cm}p{1.3cm}p{1.8cm}p{1.4cm}p{1.3cm}p{1.6cm}p{1.4cm}p{1.7cm}}
\hline
Environment & Period & Soil & AWC & Cultivar & Mean rain & Field obs. & Split \\
            &        &      & (mm) &          & (mm)      &            & \\
\hline
Castro A819 & 2006--2024 & synthetic clayey & 198 & 990005, MG 5 & 950.0 & 0 & train 2006--2016; valid 2017--2020; test 2021--2024 \\
\hline
\end{tabular}
\end{table}

DSSAT-CSM with the CROPGRO-Soybean module is the process representation. For each candidate strategy a complete treatment is generated (daily weather, soil profile, initial soil water, cultivar, sowing date, population and spacing, prescribed irrigation events, simulation controls), the season is simulated, and grain yield and seasonal irrigation are returned: $\mathbf{z} = f_{\text{DSSAT}}(y,a,\psi)$, where $y$ is the Castro season, $a$ the management action, and $\psi$ the fixed soil, cultivar, and model parameters. Simulations are executed through \texttt{DSSATTools}; the policy search uses Stable-Baselines3 and PyTorch (versions in Section~\ref{subrepro}).

Three concepts are kept distinct throughout: \emph{crop-model calibration}, the adjustment of DSSAT soil and cultivar parameters against field-scale observations, which was not possible with the available data; \emph{statistical correction}, an external transformation between simulated output and the municipal series, used in the previous diagnostic configuration and removed from the objective of the final configuration (Section~\ref{subobjective}); and \emph{policy performance}, the primary outcome, yield and water use of a management rule relative to baselines in Castro seasons not used for development.

%%============================================================
\subsection{DSSAT configuration and crop-model validation status}
\label{subcropmodel}

The executable workflow does not contain field data for crop-model calibration or independent validation. No measured Castro soil profile, cultivar-specific phenology, field or trial yield, or observed soil-water series is available. Therefore, this manuscript does not report DSSAT calibration metrics such as phenology RMSE, yield RMSE, NRMSE, bias, concordance coefficient, or soil-water error. The DSSAT outputs are simulations of a generic cultivar on a synthetic profile under the observed Castro weather; they define the response surface on which the management rules are compared, not an agronomically validated prediction of field yield.

The soil is a synthetic clayey profile created in \texttt{src/dssat\_adapter.py}, with five layers to 150 cm and 198 mm of available water between lower limit and drained upper limit; its layer parameters are exported to Supplementary Table S2. The cultivar is the generic maturity-group-5 entry \texttt{990005} from \texttt{SBGRO048.CUL}. Simulation controls are: soil-water balance on (\texttt{WATER = Y}), nitrogen dynamics off (\texttt{NITRO = N}), no fertilization, sowing prescribed (\texttt{PLANT = R}), and irrigation prescribed only when events exist (\texttt{IRRIG = R}; otherwise \texttt{IRRIG = N}). Nitrogen is therefore assumed non-limiting by configuration.

\emph{Soil-water spin-up and simulation window.} Each simulation starts on August 1 of the sowing year with every soil layer at its drained upper limit and runs through the sowing date to physiological maturity. The weather file passed to DSSAT covers August 1 of the sowing year to June 30 of the following year, so that late sowings have weather until maturity. The soil water at sowing is therefore produced by the observed August--sowing weather rather than set to a default value on the sowing day. In the previous diagnostic configuration the simulation started on the sowing date with default initial conditions and the weather file ended on April 30.

The sowing action window is September 15--December 20. The script \texttt{src/planting\_window\_audit.py} records the MAPA soybean ZARC ordinance for Paraná, safra 2024/2025 (Portaria SPA/MAPA no. 118; \url{https://www.gov.br/agricultura/pt-br/assuntos/riscos-seguro/programa-nacional-de-zoneamento-agricola-de-risco-climatico/portarias/safra-vigente/parana/PORTN118SOJAPR.ret2.pdf}), extracts the general cycle classes reported in the ordinance (100, 115 and 130 days) and the AD1--AD6 available-water classes, and checks the sowing dates against the coded window. The PDF text does not contain the municipality-specific Castro table; the ordinance directs users to the ZARC Oficial panel for municipality-, soil- and group-specific recommendations.

%===============================================================
\subsubsection{Weather data and quality}
\label{subweather}

Hourly meteorological observations are obtained from INMET automatic station A819 in Castro, Paraná, Brazil, available at \url{https://portal.inmet.gov.br/dadoshistoricos} (accessed on 26 September 2026). Annual files from July 2006 to May 2026 are consolidated; the 2022 records and the first day of the series are taken from the complementary station export \path{dados_A819_H_2006-07-08_2026-01-01.csv}, used only to fill absent timestamps. The consolidated hourly file contains 174,432 records and covers the period from 8 July 2006 00:00 to 31 May 2026 23:00. Sentinel values are converted to missing; values $\leq -90$, negative precipitation, temperatures outside $-20$ to $50\,^{\circ}$C, and radiation below zero or above 6000 kJ m$^{-2}$ h$^{-1}$ are removed. Daily precipitation and radiation are hourly sums (radiation converted to MJ m$^{-2}$); daily maximum and minimum temperature are hourly extremes. Days with fewer than 18 valid precipitation or temperature hours, or fewer than 6 valid radiation hours, are set to missing. Missing precipitation is filled with the day-of-year climatological mean (monthly mean as fallback); missing or sub-0.5 MJ m$^{-2}$ d$^{-1}$ radiation is estimated by the Hargreaves method and bounded to 1--35 MJ m$^{-2}$ d$^{-1}$; missing extremes are replaced by the daily mean and interpolated.

Filling missing precipitation with a climatological mean can remove a real dry spell, and interpolating temperature across a long gap can produce values that were never observed. The diagnostic script exports the weather-completeness table by season to \texttt{weather\_completeness\_by\_season.csv}; the precipitation-imputation sensitivity (\texttt{src/weather\_imputation\_sensitivity.py}) compares day-of-year climatology, monthly climatology, interpolation of short gaps followed by day-of-year climatology, and zero filling as a dry lower bound. The coverage of the pre-season windows used in the context is analysed separately in Section~\ref{subres_context}.

%===============================================================
\subsubsection{Municipal yield series (contextual verification only)}
\label{subsidra}

Observed soybean yields are obtained from Table 1612 of the IBGE Automatic Recovery System (SIDRA), available at \url{https://sidra.ibge.gov.br/tabela/1612} (accessed on 26 September 2026); only the municipal average soybean yield of Castro (kg ha$^{-1}$) is used. The series carries a technological trend unrelated to the weather of each season, and its unit of observation (municipal annual average over many fields, cultivars, and management levels) does not coincide with the unit of simulation (one field, one management). It is therefore used as contextual information on the range of yields, not as a target for calibrating management response, and it enters neither DSSAT nor the policy objective. A season is identified by its sowing year $y$ (sowing September--December of $y$, harvest January--April of $y+1$). Following the IBGE PAM rule that assigns production to the civil year in which most production is registered (\url{https://www.ibge.gov.br/biblioteca/visualizacao/periodicos/66/pam_2023_v50_br_notas_tecnicas.pdf}), season $y$ corresponds to the SIDRA value of year $y+1$; the audit table \texttt{season\_sidra\_pairing\_audit.csv} stores both candidate pairings.

%%============================================================
\subsection{Decision formulation and information set}
\label{subdecision}

The decision is treated as a one-step contextual problem: on the decision date $t_0$ the policy observes a context $s$, chooses an action $a$, the season is simulated, a terminal outcome $r(s,a)$ is returned, and the episode ends. No temporal credit assignment is involved. The term \emph{context-conditioned policy search} is used for the method; reinforcement learning refers to the PPO implementation used to search the stochastic policy.

\emph{Information set (context).} The context contains six variables available on $t_0$:
\begin{equation}
s = \big(P_{30},\; P_{90},\; \bar{T}_{30},\; \bar{R}_{90},\; \mathrm{ONI}_{\mathrm{MJJ}},\; \mathrm{ONI}_{\mathrm{MJJ}}-\mathrm{ONI}_{\mathrm{FMA}}\big),
\label{eq_context}
\end{equation}
where $P_{30}$ and $P_{90}$ are station rainfall accumulated over the 30 and 90 days ending on August 31 (mm), $\bar{T}_{30}$ is the mean air temperature of the 30 days ending on August 31 ($^{\circ}$C), $\bar{R}_{90}$ is the mean daily global radiation of the 90 days ending on August 31 (MJ m$^{-2}$ d$^{-1}$), and $\mathrm{ONI}_{\mathrm{MJJ}}$ is the Oceanic Ni\~no Index for May--July of the sowing year, with its change since February--April as an indicator of the ENSO tendency. The ONI series is the NOAA/CPC index based on ERSST, obtained from NOAA/PSL (\url{https://psl.noaa.gov/data/correlation/oni.data}); the May--July value is published in early August and is therefore available on $t_0$, whereas the June--August value is published after $t_0$ and is not used. Each variable is standardized with the mean and standard deviation of the eleven training seasons; the scaler is stored with each trained model and applied unchanged to validation and test contexts. The climatological rainfall and temperature terms of the previous context version were removed because they are nearly constant across seasons and, after standardization, would only transmit noise. Realized weather after $t_0$ is never part of the context.

\emph{Action.} The policy outputs four continuous values in $[-1,1]$ mapped linearly to: sowing offset $d \in \{0,\dots,96\}$ days after September 15 (rounded to the nearest day); irrigation trigger $\theta \in [10, 150]$ mm (rounded to 1 mm); depth per event $l \in [0, 35]$ mm (rounded to 0.5 mm); and seasonal cap $L \in [20, 260]$ mm (rounded to 5 mm). Events with depth below 5 mm are not applied, so that rainfed management corresponds to a region of the continuous action space rather than to a single corner; when no event is scheduled the environment passes an empty schedule to DSSAT and $W = 0$. The upper bound of $\theta$ was raised from 90 mm in the previous configuration to 150 mm, about three quarters of the available water of the profile, after the development grid showed that triggers $\leq 90$ mm irrigated in seasons without a yield response (Section~\ref{subres_landscape}). Rounding makes repeated actions map to identical DSSAT treatments, which are then read from an in-memory cache.

\emph{Irrigation rule.} The three irrigation parameters define a parameterized schedule computed from a proxy water balance on the daily weather from the sowing date, $B_{d+1} = \max(0, B_d + E_d - P_d)$ with $E_d = 0.16\,R_d + 0.08\,\max(\bar{T}_d - 10, 0)$ (mm d$^{-1}$, $R_d$ in MJ m$^{-2}$ d$^{-1}$, $\bar{T}_d$ in $^{\circ}$C), inspected every 4 days, the return interval of a centre pivot under supplemental irrigation: if $B_d \geq \theta$ and the accumulated irrigation is below $L$, an event of $\min(l, L - \sum W)$ mm is scheduled on day $d$ and $B_d$ is reduced by that depth. The decision to irrigate at the start of day $d$ uses $B_d$, computed exclusively from the weather observed up to the end of day $d-1$; no rainfall or temperature observed on or after day $d$ alters the event decided at the start of day $d$. The rule is evaluated over the 150 days after sowing, the events are passed to DSSAT as prescribed irrigation, and the rule is not a feedback on the simulated soil water. The proxy coefficients were not calibrated against reference evapotranspiration or observed soil-water deficit. In the previous configuration the inspection interval was 12 days.

%%============================================================
\subsection{Policies and baselines}
\label{subbaselines}

Comparators are separated by the information they use, so that the value of a good global rule, the value of context, and the value of the algorithm can be attributed. All comparators except PPO are evaluated on a common grid of 602 candidate rules: 13 sowing offsets (0, 8, \dots, 88, 96 days after September 15) combined with rainfed management or with 45 irrigation rules ($\theta \in \{30, 60, 90, 120, 150\}$ mm, $l \in \{15, 25, 35\}$ mm, $L \in \{80, 160, 240\}$ mm), plus the four constructed schedules. Each candidate is simulated once per season with the same weather, soil, cultivar, spin-up and objective as the PPO environment.

\emph{Without context:} (1) four \emph{constructed reference schedules}: sowing October 15 rainfed, and sowing October 15, September 25, or November 10 with the irrigation rule set to trigger 45 mm, 18 mm per event, cap 120 mm. These dates and parameters were chosen by the authors to span the action window and a moderate irrigation level; they are not agronomic recommendations. (2) An \emph{optimized fixed rule}: the grid candidate with the highest mean objective over the development seasons. For the validation comparison the selection uses only the eleven training seasons; for the test comparison it uses all fifteen development seasons (2006--2020).

\emph{With the same pre-season context:} (3) a \emph{contextual nearest-neighbour rule} ($k$NN): for a target season, the three development seasons closest in the standardized context of Eq.~(\ref{eq_context}) are identified, and the grid candidate with the highest mean objective over those three seasons is applied. The rule has one parameter ($k = 3$), fixed before the analysis, and receives exactly the information given to PPO. (4) The PPO contextual policy.

\emph{Retrospective ceiling, never a competitor:} (5) a per-season oracle, the grid candidate with the highest objective in each season, selected with the realized weather of that season. The oracle defines the regret and the gap that adaptation could, at most, close within the grid.

\begin{table}[!h]
\centering
\footnotesize
\caption{Comparative design of policies and baselines. Budgets count DSSAT simulation requests.}
\label{tab_design}
\begin{tabular}{p{3.6cm}p{3.3cm}p{2.5cm}p{3.0cm}p{1.4cm}}
\hline
Comparator & Information used & Action space & DSSAT budget & Adaptive \\
\hline
Constructed schedules (4) & none & fixed & 1 run per season & no \\
Optimized fixed rule & development seasons & grid (602 rules) & 602 per season & no \\
Contextual $k$NN rule & context on $t_0$; development seasons & grid (602 rules) & same grid & yes \\
PPO contextual policy & context on $t_0$; training seasons; validation for checkpoint & continuous, 4 dimensions & 30,780--36,936 requests per seed & yes \\
Retrospective oracle (ceiling) & realized weather of the season & grid (602 rules) & same grid & retrospective \\
\hline
\end{tabular}
\end{table}

%%============================================================
\subsection{Proximal Policy Optimization for the one-step decision}
\label{subppo}

\emph{Policy and value function.} The policy $\pi_\phi(a \mid s)$ is a diagonal Gaussian over the four pre-squashing action components, with mean $\mu_\phi(s)$ given by a multilayer perceptron with two hidden layers of 32 units and hyperbolic-tangent activation, and with a state-independent log standard deviation vector initialized at $-0.5$ (standard deviation 0.61). Actions are sampled from the Gaussian and clipped to $[-1,1]$ before decoding. A separate network of the same architecture estimates the value $V_\varphi(s)$, the expected objective of the context under the current policy.

\emph{Advantage.} Because each episode has one step and terminates after the simulation, the return of an episode is its objective $r(s,a)$ and the advantage reduces to
\begin{equation}
\hat{A}(s,a) = r(s,a) - V_\varphi(s).
\label{eq_advantage}
\end{equation}
The discount factor ($\gamma = 0.99$) and the generalized advantage estimator ($\lambda = 0.95$) of the implementation have no effect in this formulation. Advantages are normalized within each minibatch.

\emph{Objective.} With $\rho(\phi) = \pi_\phi(a \mid s)/\pi_{\phi_{\text{old}}}(a \mid s)$ the probability ratio between the updated policy and the policy that collected the data, PPO maximizes the clipped surrogate
\begin{equation}
\mathcal{L}(\phi,\varphi) = \mathbb{E}\Big[\min\big(\rho(\phi)\hat{A},\; \mathrm{clip}(\rho(\phi), 1-\epsilon, 1+\epsilon)\,\hat{A}\big)\Big] - c_v\, \mathbb{E}\Big[\big(V_\varphi(s) - r\big)^2\Big] + c_e\, \mathbb{E}\big[\mathcal{H}(\pi_\phi(\cdot \mid s))\big],
\label{eq_ppo}
\end{equation}
with $\epsilon = 0.2$, $c_v = 0.5$ and $c_e = 0$. Each update collects 512 episodes (8 parallel environments $\times$ 64 episodes), performs 10 epochs of minibatch gradient steps with minibatches of 128 episodes, uses Adam with learning rate $3\times 10^{-4}$, and clips the gradient norm at 0.5.

\emph{Training data and regularization.} Each training episode draws one of the eleven training seasons at random, so the network sees eleven distinct contexts. To limit the memorization of season-specific actions, Gaussian noise with standard deviation 0.75 is added to the standardized context of every training episode; validation and test contexts are not perturbed. The noise level and the network size were chosen on the validation seasons in one development comparison with seed 42 (noise 0.25 and two layers of 64 units versus noise 0.75 and two layers of 32 units; Section~\ref{subres_learning}); no other hyperparameter was tuned.

\emph{Validation, checkpoint and stopping.} Every 2,048 training steps the deterministic policy (the Gaussian mean) is evaluated once on each of the four validation seasons; DSSAT is deterministic, so repeated evaluations of the same season are not drawn. The checkpoint with the highest mean validation objective is retained. Training runs for at least 20,000 steps and stops after 12 consecutive evaluations without improvement or at 80,000 steps.

\emph{Parallel simulation.} The eight training environments run in separate processes, each with its own DSSAT working directory and its own copy of the DSSAT data folder, because \texttt{DSSATTools} rebuilds that folder at import time and concurrent processes would otherwise overwrite each other's files. Five seeds (42--46) were trained sequentially with otherwise identical configurations.

%%============================================================
\subsection{Objective, water constraint, and risk metric}
\label{subobjective}

\emph{Outcome hierarchy.} \emph{Primary results:} yield (kg ha$^{-1}$), seasonal irrigation (mm), and their paired differences between the contextual policy and each comparator on the frozen chronological test block (Eq.~\ref{eq_paired}). \emph{Systemic interpretation:} the productivity--water relation, the regret with respect to the retrospective oracle, and the behaviour of the policy across contrasting seasons. \emph{Secondary operational result:} the scalar $r$ of Eq.~(\ref{eq_reward}), the optimization function of the search. \emph{Stability outcomes:} DSSAT failures, action saturation, and variation between seeds.

\emph{Optimization function.} The scalar optimized by the search for one season is
\begin{equation}
r = \frac{Y}{Y_{\text{ref}}} - \kappa\, W,
\label{eq_reward}
\end{equation}
with $Y$ the raw DSSAT grain yield (kg ha$^{-1}$), $Y_{\text{ref}} = 4200$ kg ha$^{-1}$ a normalization constant, $W$ the seasonal irrigation (mm), and $\kappa = 0.0015$ mm$^{-1}$, the value registered before the final run. With these values one millimetre of irrigation must raise yield by at least 6.3 kg ha$^{-1}$ to increase $r$. A DSSAT run that fails to complete receives $r = -2$; no failure occurred in the final evaluations. The water-penalty sensitivity re-optimizes the fixed rule on the development seasons for $\kappa \in \{0, 0.0005, 0.0015, 0.003, 0.006\}$ mm$^{-1}$ and re-scores the PPO points without retraining.

\emph{Removal of the statistical correction.} In the previous configuration $Y$ was replaced by a ridge regression of the municipal yield on thirteen predictors, including year trend, sowing day and irrigation, fitted on training seasons at five fixed sowing dates. Because the five rows of each season shared one municipal target, the regression carried no information on the response to sowing date or irrigation; its irrigation coefficients were negative ($-0.106$ per mm; rain--irrigation interaction $-0.361$), and within a season it reduced the spread across sowing dates from up to 1,578 kg ha$^{-1}$ of raw yield to 2--48 kg ha$^{-1}$ of corrected yield. The final configuration uses raw DSSAT yield in Eq.~(\ref{eq_reward}); a constant scale or bias correction would not change the ranking of actions within a season.

\emph{Regret and risk.} Regret is the difference between the objective (or yield) of the retrospective oracle and that of the policy in the same season. The risk metrics are the 10th percentile of yield and CVaR$_{10}$ (mean of the yields at or below the 10th percentile), computed over seasons (and seeds, for PPO) in the test block.

%%============================================================
\subsection{Experimental design, temporal split, and uncertainty}
\label{subdesign}

\emph{Unit of analysis.} The independent agronomic unit is the Castro season. A PPO episode, a validation evaluation, or a simulated row of the same season is not an independent unit; a training seed is a replication of the optimization procedure applied to the same contexts and does not increase the agronomic sample size.

\emph{Temporal split.} Seasons are split chronologically according to the registry written by \texttt{src/temporal\_split\_registry.py} (Table~\ref{tab_temporal_registry_castro}): training 2006--2016 (11 seasons), validation 2017--2020 (4), and test 2021--2024 (4). The 2025 season has no observed municipal yield and is not used. Because the study has data for Castro only, no spatial split is part of the design.

\begin{table}[!h]
\centering
\scriptsize
\caption{Temporal registry for the Castro run.}
\label{tab_temporal_registry_castro}
\begin{tabular}{llrrr}
\hline
Block & Years & $n$ & Mean rainfall & Observed yield years \\
      &       &     & (mm season$^{-1}$) & \\
\hline
Development train & 2006--2016 & 11 & 892.75 & 11/11 \\
Development validation & 2017--2020 & 4 & 1024.75 & 4/4 \\
Test & 2021--2024 & 4 & 980.10 & 4/4 \\
\hline
\end{tabular}
\end{table}

\emph{Governance of the test block.} In an earlier version of the workflow the 2021--2024 block was evaluated once with a policy that received realized-season weather and an objective based on the statistical correction. That evaluation identified no tuning target in the block; the changes introduced afterwards (raw yield in the objective, the pre-season context of Eq.~\ref{eq_context}, the spin-up and weather window, the 4-day inspection, the trigger range, and the regularization of PPO) were decided on the development seasons, using the development grid and the validation curves. For the configuration reported here the test block was opened once, after the five seeds had been trained and the comparators selected. The context-coverage analysis of Section~\ref{subres_context} was defined after that opening and is reported as a post hoc sensitivity.

\emph{Paired effects.} For each season $y$, seed $s$, metric $M$, and comparator $b$, the paired effect is
\begin{equation}
\Delta M_{y,s,b} = M_{\pi,y,s} - M_{b,y},
\label{eq_paired}
\end{equation}
where positive values indicate higher yield, objective, or irrigation for PPO.

\emph{Decomposition of the paired effects.} For each comparator the paired yield differences are decomposed as
\begin{equation}
\Delta_{y,s} = \mu + v_y + w_s + \varepsilon_{y,s},
\label{eq_mixed}
\end{equation}
where $\mu$ is the mean paired difference, $v_y$ the season effect (season mean minus $\mu$), $w_s$ the seed effect (seed mean minus $\mu$), and $\varepsilon_{y,s}$ the residual. With four seasons, the components are reported as standard deviations and all paired effects are shown, without interval inference.

\emph{Value decomposition.} The total gain over the agronomic calendar (rainfed October 15) is decomposed as the gain from choosing a good fixed rule (optimized fixed rule minus calendar), the gain from context (contextual $k$NN minus optimized fixed rule), and the incremental value of the algorithm (PPO minus contextual $k$NN). The share of the oracle gain captured by a rule is its gain over the calendar divided by the gain of the oracle over the calendar.

\emph{Skill of the context.} Before interpreting the test results, the association between each context variable and three season outcomes of the development grid (best rainfed sowing offset, irrigation gain of the best irrigated rule over the best rainfed rule, and yield of the rainfed October 15 schedule) is summarized by Spearman correlations over the fifteen development seasons. The position of every season's context relative to the training seasons is reported as the number of variables outside the training range and the Euclidean distance to the nearest training season in standardized units.

%%============================================================
\subsection{Reproducibility and configuration}
\label{subrepro}

The configuration \texttt{configs/experiment\_castro\_t0.yaml} records the decision date, context variables, action bounds, irrigation-rule parameters, spin-up and weather window, objective, water penalty, PPO hyperparameters, seeds, failure penalty, and DSSAT settings. The script \texttt{run\_final\_castro\_t0.ps1} trains the five seeds (\texttt{src/run\_multiseed.py}), collects their evaluations (\texttt{src/summarize\_multiseed.py}), and evaluates the comparators and paired effects on the test block (\texttt{src/final\_comparators.py}); \texttt{src/article\_results.py} produces the tables and figures, and \texttt{src/context\_qc\_sensitivity.py} the post hoc context analysis. Per-seed configurations are generated in \texttt{configs/\_generated\_ppo\_castro\_t0}.

The DSSAT log records DSSAT-CSM Ver. 4.8.0.028; the Python stack includes \texttt{DSSATTools} 3.0.2, Stable-Baselines3 2.7.1, PyTorch 2.10.0, Gymnasium 1.2.2, pandas 2.3.3, NumPy 1.26.4, Matplotlib 3.10.9, and TensorBoard 2.20.0. Each seed took 48--56 min on a 12-core workstation with eight simulation processes. The repository includes \texttt{environment.yml}, an Apache~2.0 license file, the generated configurations, the validation curves and TensorBoard logs of every seed, the candidate grids, and the comparator tables.

\begin{figure}[!h]
\centering
\resizebox{\linewidth}{!}{%
\begin{tikzpicture}[
    font=\scriptsize,
    box/.style={rectangle, rounded corners=1pt, draw=black, line width=0.4pt, align=center, text width=2.7cm, minimum height=0.6cm, inner sep=2pt},
    data/.style={box, fill=green!8},
    proc/.style={box, fill=blue!7},
    dec/.style={box, fill=red!7},
    outp/.style={box, fill=gray!10},
    arrow/.style={-Latex, line width=0.45pt},
    darrow/.style={-Latex, line width=0.45pt, dashed}
]

\node[data] (weather) at (0,0) {INMET hourly weather and NOAA ONI},
\node[proc] (qc) at (0,-1.35) {Quality control, daily series, imputation},
\node[data] (field) at (4.4,0) {Synthetic soil profile and generic cultivar 990005},
\node[proc] (calib) at (4.4,-1.35) {DSSAT-CROPGRO configuration (not field-calibrated)},
\node[data] (sidra) at (8.8,0) {SIDRA municipal yield},
\node[outp] (verif) at (8.8,-1.35) {Contextual information only (not an input to DSSAT or to the policy)},
\draw[arrow] (weather) -- (qc),
\draw[arrow] (field) -- (calib),
\draw[darrow] (sidra) -- (verif),
% --- decision timeline ---
\node[proc] (init) at (0,-3.1) {August 1: soil-water spin-up from drained upper limit},
\node[dec] (t0) at (4.4,-3.1) {Decision date $t_0$ = Sep 1: standardized context $s$ (Eq.~\ref{eq_context})},
\node[dec] (policy) at (8.8,-3.1) {PPO policy or comparator $a=\pi(s)$},
\node[dec] (action) at (13.2,-3.1) {Action $[d,\theta,l,L]$; $l<5$ mm = rainfed},
\node[proc] (dssat) at (13.2,-4.85) {DSSAT season with causal irrigation schedule (4-day inspection)},
\node[outp] (out) at (8.8,-4.85) {Outcomes: raw yield, irrigation, $r$, risk},
\node[outp] (eval) at (4.4,-4.85) {Paired evaluation on the 2021--2024 test block},
\draw[arrow] (init) -- (t0),
\draw[arrow] (t0) -- (policy),
\draw[arrow] (policy) -- (action),
\draw[arrow] (action) -- (dssat),
\draw[arrow] (dssat) -- (out),
\draw[arrow] (out) -- (eval),
\draw[darrow] (out.north) -- node[right,font=\tiny]{PPO update (training seasons); checkpoint (validation)} (policy.south),
\draw[arrow] (qc.west) -- ++(-0.35,0) |- ([yshift=-0.75cm]dssat.south) -- (dssat.south),
\draw[arrow] (calib.south) -- ++(0,-0.25) -| ([xshift=0.9cm]dssat.south) ,
\draw[line width=0.6pt] (-1.5,-2.35) -- (14.7,-2.35),
\node[font=\tiny, anchor=south west] at (-1.5,-2.33) {information available before $t_0$},
\node[font=\tiny, anchor=south east] at (14.7,-2.33) {information observed during the season (simulator and causal irrigation rule)},
\end{tikzpicture}}
\caption{System and decision timeline of the framework. Hourly weather passes through quality control to the context and to DSSAT; soil and cultivar are generic and not calibrated with field data; the SIDRA municipal series enters neither DSSAT nor the policy. The decision is taken on $t_0$ from the standardized pre-season context, the season is simulated after a soil-water spin-up from August 1, the irrigation rule is causal in time, PPO is updated on the training seasons and its checkpoint is selected on the validation seasons, and the paired evaluation uses the 2021--2024 test block.}
\label{fig:flowchart_dssat_ppo}
\end{figure}


% =====================================================================
\section{Results}
\label{resultados}

The results are presented in the order in which the evidence was produced: the response surface of the development seasons, the learning behaviour of PPO, performance on the training and validation seasons, the paired evaluation on the test block, and the analysis of the pre-season context that explains the test results.

%%============================================================
\subsection{Response surface of the development seasons}
\label{subres_landscape}

With raw DSSAT yield, the sowing date and irrigation produced large differences within a season (Table~\ref{tab_landscape}, Fig.~\ref{fig_rainfed_response}). Across the fifteen development seasons, the range of rainfed yield across the thirteen sowing offsets was 1,853 kg ha$^{-1}$ on average (880--3,932 kg ha$^{-1}$). Irrigation increased the best attainable yield by more than 100 kg ha$^{-1}$ in three seasons only: 2007 (+2,820 kg ha$^{-1}$), 2013 (+2,432 kg ha$^{-1}$) and 2014 (+466 kg ha$^{-1}$). In the other twelve seasons the best rainfed and the best irrigated rule differed by 0--62 kg ha$^{-1}$. The previous statistical correction had compressed the within-season spread to 2--48 kg ha$^{-1}$, which left the search without a signal to separate sowing dates or irrigation rules.

\begin{table}[!h]
\centering
\scriptsize
\caption{Raw DSSAT response in the development seasons (grid of 602 rules). Sowing range: maximum minus minimum rainfed yield across the 13 sowing offsets; irrigation gain: best grid rule minus best rainfed rule.}
\label{tab_landscape}
\begin{tabular}{lrrrr}
\hline
Season & Sowing range & Best rainfed & Best rule & Irrigation gain \\
       & (kg ha$^{-1}$) & (kg ha$^{-1}$) & (kg ha$^{-1}$) & (kg ha$^{-1}$) \\
\hline
2006 & 3,932 & 5,367 & 5,367 & 0 \\
2007 & 1,887 & 2,795 & 5,615 & 2,820 \\
2008 & 880 & 5,280 & 5,300 & 20 \\
2009 & 3,924 & 5,038 & 5,100 & 62 \\
2010 & 1,765 & 5,428 & 5,428 & 0 \\
2011 & 1,112 & 6,004 & 6,004 & 0 \\
2012 & 1,051 & 4,228 & 4,228 & 0 \\
2013 & 1,372 & 3,659 & 6,091 & 2,432 \\
2014 & 2,126 & 5,178 & 5,644 & 466 \\
2015 & 1,264 & 4,969 & 4,969 & 0 \\
2016 & 1,398 & 5,381 & 5,381 & 0 \\
2017 & 1,665 & 5,089 & 5,100 & 11 \\
2018 & 2,158 & 5,449 & 5,449 & 0 \\
2019 & 2,036 & 5,200 & 5,200 & 0 \\
2020 & 1,226 & 5,194 & 5,194 & 0 \\
\hline
\end{tabular}
\end{table}

\begin{figure}[!h]
\centering
\includegraphics[width=0.62\linewidth]{fig_rainfed_sowing_response.png}
\caption{Rainfed DSSAT yield as a function of the sowing offset after September 15 in the development seasons. Thin lines are seasons; thick lines are the means of the training (2006--2016) and validation (2017--2020) seasons.}
\label{fig_rainfed_response}
\end{figure}

The mean response to sowing date differed between the training and validation blocks. In the training seasons the mean rainfed objective decreased monotonically from the first sowing offset (1.069 at September 15 to 0.773 at December 20); in the validation seasons it was highest at offsets of 16--32 days (1.217--1.224, October 1--17) and lower at September 15 (1.175). The optimized fixed rule selected on the training seasons alone (September 15 sowing, $\theta = 120$ mm, $l = 25$ mm, $L = 240$ mm) therefore reached 1.130 on the training seasons and 1.175 on validation, below the rainfed October 15 schedule (1.221). On the training seasons the triggers $\leq 90$ mm applied water in years without yield response; with $\theta \geq 110$ mm irrigation was concentrated in 2007, 2009 and 2013 and was zero in all validation seasons, which motivated the extension of the trigger range.

%%============================================================
\subsection{Learning behaviour and checkpoint selection}
\label{subres_learning}

Figure~\ref{fig_learning_curves} shows the learning curves of the five seeds. The mean objective of the training episodes, which includes the exploration noise of the Gaussian policy and the context noise, rose from 0.93--0.99 in the first update to 1.13--1.16 at the end of training. The deterministic validation objective rose to its maximum within the first 6,144--12,288 steps and then declined towards the level of the optimized fixed rule (1.16--1.19) while the training objective kept increasing. Early stopping ended training at 30,720--36,864 steps (Table~\ref{tab_checkpoints}). The selected checkpoints reached validation objectives of 1.221--1.229, with mean validation yields of 5,130--5,162 kg ha$^{-1}$ and no irrigation; the between-seed spread of the selected validation objective was 0.008.

\begin{figure}[!h]
\centering
\includegraphics[width=\linewidth]{fig_learning_curves.png}
\caption{Learning curves of PPO for seeds 42--46. (a) Mean objective of training episodes (stochastic actions, perturbed context) recorded by TensorBoard after each update of 512 episodes. (b) Objective of the deterministic policy on the four validation seasons, evaluated every 2,048 steps; stars mark the selected checkpoints. Horizontal lines: validation objective of the retrospective oracle, the rainfed October 15 schedule, and the fixed rule optimized on the training seasons.}
\label{fig_learning_curves}
\end{figure}

\begin{table}[!h]
\centering
\scriptsize
\caption{Selected checkpoints and training budget per seed. Simulation requests count training episodes plus validation evaluations; repeated actions are read from the cache and do not call DSSAT again.}
\label{tab_checkpoints}
\begin{tabular}{rrrrrrr}
\hline
Seed & Selected step & Valid. objective & Valid. yield & Valid. irrigation & Stop step & Simulation requests \\
     &               &                  & (kg ha$^{-1}$) & (mm) & & \\
\hline
42 & 12,288 & 1.228 & 5,158 & 0.0 & 36,864 & 36,936 \\
43 & 6,144 & 1.226 & 5,149 & 0.0 & 30,720 & 30,780 \\
44 & 8,192 & 1.229 & 5,162 & 0.0 & 32,768 & 32,832 \\
45 & 6,144 & 1.227 & 5,154 & 0.0 & 30,720 & 30,780 \\
46 & 6,144 & 1.221 & 5,130 & 0.0 & 30,720 & 30,780 \\
\hline
\end{tabular}
\end{table}

The optimization diagnostics were similar across seeds (Fig.~\ref{fig_ppo_diagnostics}). The standard deviation of the Gaussian policy decreased from 0.61 to 0.44--0.47, the entropy term increased accordingly, the approximate Kullback--Leibler divergence between successive policies remained between 0.004 and 0.027, and the fraction of clipped ratios increased from 0.01--0.04 to 0.07--0.13. The explained variance of the value network rose from negative values in the first updates to 0.31--0.45, indicating that the critic captured part of the season-to-season variation of the objective.

\begin{figure}[!h]
\centering
\includegraphics[width=\linewidth]{fig_ppo_diagnostics.png}
\caption{PPO optimization diagnostics recorded by TensorBoard for seeds 42--46: standard deviation of the Gaussian policy, entropy loss, approximate Kullback--Leibler divergence between successive policies, and explained variance of the value network.}
\label{fig_ppo_diagnostics}
\end{figure}

The regularization was decided in one development comparison with seed 42 (Fig.~\ref{fig_regularization}). With context noise 0.25 and two hidden layers of 64 units, the validation objective stayed between 1.145 and 1.178 and the best checkpoint was the second evaluation (2,048 steps). The final policy of that run reached 1.187 on the training seasons, close to the oracle (1.213), but 1.164 on validation: it irrigated 245 mm in 2007 and 160 mm in 2013, the two drought seasons of the training block, and sowed on September 15--22 in seven of the eleven training seasons and in three of the four validation seasons. With context noise 0.75 and two layers of 32 units, the validation objective reached 1.228.

\begin{figure}[!h]
\centering
\includegraphics[width=0.62\linewidth]{fig_regularization_seed42.png}
\caption{Validation objective of seed 42 with the development configuration (context noise 0.25, two hidden layers of 64 units) and with the final configuration (context noise 0.75, two hidden layers of 32 units).}
\label{fig_regularization}
\end{figure}

%%============================================================
\subsection{Performance on training and validation seasons}
\label{subres_development}

Table~\ref{tab_by_split} reports the objective, yield and irrigation of the selected PPO checkpoints and of the comparators in each block. On the training seasons PPO reached a mean objective of 1.156 (1.150--1.162 across seeds), above the optimized fixed rule (1.130) and the rainfed October 15 schedule (1.020) and below the oracle (1.213). The selected checkpoints applied 150--260 mm in 2007 and 2013, 25--72 mm in 2014 and 2015, and less than 40 mm in the other seasons, and they sowed between September 15 and November 16. On the validation seasons PPO reached 1.226 (1.221--1.229), above the rainfed October 15 schedule (1.221), the contextual $k$NN rule (1.181) and the fixed rule optimized on the training seasons (1.175), with no irrigation and sowing between September 15 and October 13.

\begin{table}[!h]
\centering
\scriptsize
\caption{Mean objective $r$, yield and irrigation by block. Training and validation comparators are selected on the training seasons; test comparators on all development seasons (2006--2020). PPO values are means over five seeds (range of seed means in parentheses).}
\label{tab_by_split}
\begin{tabular}{llrrr}
\hline
Block & Policy & $r$ & Yield (kg ha$^{-1}$) & Irrigation (mm) \\
\hline
Training (11) & PPO & 1.156 (1.150--1.162) & 5,150 & 47.0 \\
              & Optimized fixed rule & 1.130 & 5,040 & 46.8 \\
              & Rainfed Oct 15 & 1.020 & 4,284 & 0.0 \\
              & Retrospective oracle & 1.213 & 5,325 & 36.4 \\
Validation (4) & PPO & 1.226 (1.221--1.229) & 5,150 & 0.0 \\
              & Optimized fixed rule & 1.175 & 4,935 & 0.0 \\
              & Contextual $k$NN rule & 1.181 & 4,960 & 0.0 \\
              & Rainfed Oct 15 & 1.221 & 5,127 & 0.0 \\
              & Retrospective oracle & 1.246 & 5,233 & 0.0 \\
Test (4)      & PPO & 0.961 (0.923--0.998) & 4,061 & 3.9 \\
              & Optimized fixed rule & 1.134 & 4,904 & 22.5 \\
              & Contextual $k$NN rule & 0.995 & 4,179 & 0.0 \\
              & Rainfed Oct 15 & 1.033 & 4,340 & 0.0 \\
              & Retrospective oracle & 1.184 & 5,137 & 26.3 \\
\hline
\end{tabular}
\end{table}

%%============================================================
\subsection{Paired evaluation on the test block}
\label{subres_test}

On the 2021--2024 test block the ordering observed on the development seasons was reversed (Table~\ref{tab_test_summary}). The optimized fixed rule selected on the fifteen development seasons (October 9 sowing, $\theta = 90$ mm, $l = 15$ mm, $L = 240$ mm) reached a mean objective of 1.134 and a mean yield of 4,904 kg ha$^{-1}$ with 22.5 mm of irrigation. The irrigated October 15 schedule reached the same objective (1.134) with higher yield and more than twice the water (5,092 kg ha$^{-1}$, 52.5 mm). PPO reached 0.961 (between-seed SD 0.035) and 4,061 kg ha$^{-1}$ (between-seed SD 168 kg ha$^{-1}$) with 3.9 mm, below the rainfed October 15 schedule (1.033; 4,340 kg ha$^{-1}$) and the contextual $k$NN rule (0.995; 4,179 kg ha$^{-1}$). PPO also had the lowest 10th percentile (2,681 kg ha$^{-1}$) and CVaR$_{10}$ (2,563 kg ha$^{-1}$) of yield.

\begin{table}[!h]
\centering
\scriptsize
\caption{Test block 2021--2024. PPO statistics are computed over 20 season--seed pairs; comparator statistics over four seasons. P10: 10th percentile of yield; CVaR$_{10}$: mean of yields at or below P10.}
\label{tab_test_summary}
\begin{tabular}{lrrrrr}
\hline
Policy & $r$ & Yield & P10 & CVaR$_{10}$ & Irrigation \\
       &     & (kg ha$^{-1}$) & (kg ha$^{-1}$) & (kg ha$^{-1}$) & (mm) \\
\hline
Retrospective oracle & 1.184 & 5,137 & 3,662 & 3,117 & 26.3 \\
Optimized fixed rule & 1.134 & 4,904 & 3,633 & 3,086 & 22.5 \\
Irrigated Oct 15 (45/18/120) & 1.134 & 5,092 & 3,680 & 3,139 & 52.5 \\
Irrigated Sep 25 (45/18/120) & 1.072 & 4,703 & 3,589 & 3,035 & 31.5 \\
Irrigated Nov 10 (45/18/120) & 1.052 & 4,795 & 3,482 & 3,003 & 60.0 \\
Rainfed Oct 15 & 1.033 & 4,340 & 3,292 & 3,139 & 0.0 \\
Contextual $k$NN rule & 0.995 & 4,179 & 3,283 & 3,117 & 0.0 \\
PPO (5 seeds) & 0.961 & 4,061 & 2,681 & 2,563 & 3.9 \\
\hline
\end{tabular}
\end{table}

The paired differences (Table~\ref{tab_paired}, Fig.~\ref{fig_test_paired}) were negative against the optimized fixed rule in all four seasons ($-843$ kg ha$^{-1}$ on average; 3 of 20 season--seed pairs with a higher objective for PPO) and against the oracle in all season--seed pairs ($-1,076$ kg ha$^{-1}$). Against the rainfed October 15 schedule and the contextual $k$NN rule the mean differences were $-279$ and $-118$ kg ha$^{-1}$, with PPO ahead in one and two seasons, respectively. The season effect dominated the decomposition of Eq.~(\ref{eq_mixed}): for the comparison with the optimized fixed rule the standard deviation of the season effects was 864 kg ha$^{-1}$, that of the seed effects 168 kg ha$^{-1}$, and that of the residuals 274 kg ha$^{-1}$.

\begin{table}[!h]
\centering
\scriptsize
\caption{Paired effects of PPO on the test block (PPO minus comparator; 4 seasons $\times$ 5 seeds). SD components follow Eq.~(\ref{eq_mixed}).}
\label{tab_paired}
\begin{tabular}{lrrrrrrr}
\hline
Comparator & $\Delta r$ & $\Delta$ yield & $\Delta$ irrigation & SD season & SD seed & SD residual & Seasons with \\
           &            & (kg ha$^{-1}$) & (mm) & (kg ha$^{-1}$) & (kg ha$^{-1}$) & (kg ha$^{-1}$) & $\Delta r > 0$ \\
\hline
Optimized fixed rule & $-0.173$ & $-843$ & $-18.6$ & 864 & 168 & 274 & 0/4 \\
Contextual $k$NN rule & $-0.034$ & $-118$ & 3.9 & 501 & 168 & 274 & 2/4 \\
Rainfed Oct 15 & $-0.072$ & $-279$ & 3.9 & 496 & 168 & 274 & 1/4 \\
Retrospective oracle & $-0.223$ & $-1,076$ & $-22.4$ & 1,027 & 168 & 274 & 0/4 \\
\hline
\end{tabular}
\end{table}

\begin{figure}[!h]
\centering
\includegraphics[width=0.85\linewidth]{fig_test_paired_yield.png}
\caption{Paired yield differences between PPO and each comparator in the test seasons. Points are seeds; horizontal bars are season means.}
\label{fig_test_paired}
\end{figure}

The season-level results identify where the differences arose (Table~\ref{tab_test_seasons}, Fig.~\ref{fig_test_sowing}). In 2021 and 2022 all PPO seeds sowed between November 15 and December 3, while the fixed rule sowed on October 9 and the oracle on October 17; PPO lost 630 and 413 kg ha$^{-1}$ relative to the fixed rule. In 2023 PPO sowed between September 18 and October 9 and was 215 kg ha$^{-1}$ below the fixed rule. In 2024, the only test season with a water deficit, rainfed management at October 15 produced 3,649 kg ha$^{-1}$, the fixed rule applied 75 mm and produced 6,213 kg ha$^{-1}$, and the oracle applied 105 mm and produced 6,709 kg ha$^{-1}$. PPO sowed between September 15 and 26 and applied 0--31 mm (mean 10.1 mm), with a mean yield of 4,099 kg ha$^{-1}$; the two seeds that applied water (19.5 and 31 mm) produced 4,463 and 4,921 kg ha$^{-1}$, and the three that did not produced 3,670--3,737 kg ha$^{-1}$.

\begin{table}[!h]
\centering
\scriptsize
\caption{Season-level yield (kg ha$^{-1}$) and irrigation (mm, in parentheses) on the test block. PPO values are means over five seeds.}
\label{tab_test_seasons}
\begin{tabular}{lrrrr}
\hline
Policy & 2021 & 2022 & 2023 & 2024 \\
\hline
PPO & 4,280 (0.0) & 2,673 (0.0) & 5,192 (5.5) & 4,099 (10.1) \\
Optimized fixed rule & 4,910 (0) & 3,086 (0) & 5,407 (15) & 6,213 (75) \\
Contextual $k$NN rule & 4,910 (0) & 3,117 (0) & 5,018 (0) & 3,670 (0) \\
Rainfed Oct 15 & 4,942 (0) & 3,139 (0) & 5,630 (0) & 3,649 (0) \\
Retrospective oracle & 4,932 (0) & 3,117 (0) & 5,789 (0) & 6,709 (105) \\
\hline
\end{tabular}
\end{table}

\begin{figure}[!h]
\centering
\includegraphics[width=0.72\linewidth]{fig_test_sowing_dates.png}
\caption{Sowing dates selected in the test seasons by PPO (one point per seed), the optimized fixed rule, the contextual $k$NN rule and the retrospective oracle.}
\label{fig_test_sowing}
\end{figure}

\emph{Value decomposition.} Relative to the rainfed October 15 calendar, choosing the optimized fixed rule added 0.101 to the objective and 564 kg ha$^{-1}$ to yield, and captured 67\% of the gap between the calendar and the oracle. Adding context through the $k$NN rule changed the objective by $-0.139$ ($-725$ kg ha$^{-1}$) relative to the fixed rule, and replacing the $k$NN rule by PPO changed it by a further $-0.034$ ($-118$ kg ha$^{-1}$). The remaining regret of PPO with respect to the oracle was 0.223 (1,076 kg ha$^{-1}$), concentrated in 2024 (0.479) and, to a lesser extent, in 2021 (0.155) and 2023 (0.150).

%%============================================================
\subsection{Productivity--water relation and water-penalty sensitivity}
\label{subres_tradeoff}

Figure~\ref{fig_yield_water} places the test-block means of every grid rule, the constructed schedules, the comparators and the five PPO seeds in the yield--irrigation plane. The optimized fixed rule lay close to the upper envelope of the grid rules at about 20--25 mm, where the envelope rises steeply; beyond 40 mm the envelope was nearly flat (about 5,000--5,100 kg ha$^{-1}$). The PPO seeds lay at 0--11 mm and 3,875--4,257 kg ha$^{-1}$, below the rainfed October 15 point.

\begin{figure}[!h]
\centering
\includegraphics[width=0.68\linewidth]{fig_test_yield_water.png}
\caption{Mean yield and mean seasonal irrigation on the test block for the 602 grid rules (grey), the constructed schedules, the comparators and the five PPO seeds. The grid points are means of rules over the test seasons and describe the attainable set; they are not available to any operational comparator.}
\label{fig_yield_water}
\end{figure}

When the fixed rule was re-optimized on the development seasons for each water penalty (Table~\ref{tab_kappa}), its test irrigation decreased from 50 mm at $\kappa = 0$ to 25 and 22.5 mm at $\kappa = 0.0005$ and 0.0015 mm$^{-1}$, and to zero for $\kappa \geq 0.003$ mm$^{-1}$, where its test yield fell to 4,120 kg ha$^{-1}$. The oracle kept 18.8--55 mm over the same range. The PPO policies, trained at $\kappa = 0.0015$ mm$^{-1}$, applied 3.9 mm, less water than the fixed rule at every $\kappa$ below 0.003 mm$^{-1}$, so their position relative to the fixed rule did not depend on the penalty within the evaluated range.

\begin{table}[!h]
\centering
\scriptsize
\caption{Water-penalty sensitivity on the test block. The fixed rule is re-optimized on the development seasons for each $\kappa$; PPO (trained at $\kappa = 0.0015$ mm$^{-1}$) is re-scored without retraining.}
\label{tab_kappa}
\begin{tabular}{rlrrr}
\hline
$\kappa$ (mm$^{-1}$) & Policy & $r$ & Yield (kg ha$^{-1}$) & Irrigation (mm) \\
\hline
0      & Optimized fixed rule (Oct 17, 60/25/240) & 1.213 & 5,095 & 50.0 \\
       & Retrospective oracle & 1.229 & 5,161 & 55.0 \\
       & PPO & 0.967 & 4,061 & 3.9 \\
0.0005 & Optimized fixed rule (Oct 9, 90/25/240) & 1.155 & 4,903 & 25.0 \\
       & Retrospective oracle & 1.211 & 5,159 & 35.0 \\
       & PPO & 0.965 & 4,061 & 3.9 \\
0.0015 & Optimized fixed rule (Oct 9, 90/15/240) & 1.134 & 4,904 & 22.5 \\
       & Retrospective oracle & 1.184 & 5,137 & 26.3 \\
       & PPO & 0.961 & 4,061 & 3.9 \\
0.003  & Optimized fixed rule (Sep 15, 150/25/80) & 0.981 & 4,120 & 0.0 \\
       & Retrospective oracle & 1.149 & 5,110 & 22.5 \\
       & PPO & 0.955 & 4,061 & 3.9 \\
0.006  & Optimized fixed rule (Sep 15, rainfed) & 0.981 & 4,120 & 0.0 \\
       & Retrospective oracle & 1.083 & 5,021 & 18.8 \\
       & PPO & 0.944 & 4,061 & 3.9 \\
\hline
\end{tabular}
\end{table}

%%============================================================
\subsection{Pre-season context: skill, coverage and post hoc sensitivity}
\label{subres_context}

\emph{Skill on the development seasons.} The association between the context variables and the season outcomes was weak (Spearman correlations over 15 seasons). The largest correlation was between 90-day pre-season rainfall and the best rainfed sowing offset ($\rho = -0.61$): seasons with wetter June--August favoured earlier sowing. The 30-day rainfall and temperature were associated with the yield of the rainfed October 15 schedule ($\rho = 0.44$ and $-0.43$). No context variable was associated with the irrigation gain ($|\rho| \leq 0.24$), and the two ONI terms had $|\rho| \leq 0.22$ for all three outcomes. The pre-season information therefore carried some signal for the sowing date and almost none for the occurrence of a drought requiring irrigation.

\emph{Position of the test contexts.} The station recorded no valid hourly temperature in August 2021 and August 2022 (0 of 720 hours in each month). After interpolation the 30-day mean temperature took values of 20.1 and 17.7~$^{\circ}$C, against 13.1--15.6~$^{\circ}$C in the fifteen development seasons, which correspond to 7.5 and 4.3 standard deviations of the training seasons. The distance of these contexts to the nearest training season was 6.45 and 4.09 standardized units, against 1.38--3.55 among the training seasons. In both seasons all seeds sowed between November 15 and December 3, the latest dates selected by any policy in any block. The 90-day mean radiation of the validation seasons 2017--2019 was 4.6--6.1 MJ m$^{-2}$ d$^{-1}$, 4.2--5.2 standard deviations below the training mean, which is consistent with a degradation of the radiation sensor rather than with the regional climate. The contexts of 2023 and 2024 were inside the training range (distance 2.11 and 2.03).

\emph{Post hoc sensitivity.} To separate the effect of the missing records from that of the policy, a feature of the context was treated as unknown, and set to the training mean, when fewer than 75\% of the days of its window were valid under the daily rule of Section~\ref{subweather}. Applied to the test contexts, this rule replaced the four station variables in 2021, 2022 and 2023 and none in 2024. The five trained policies were re-evaluated on the corrected contexts without retraining, and the $k$NN rule was recomputed with the same rule. PPO then sowed between September 15 and October 20 in all seasons, its mean objective rose from 0.961 to 1.008 and its mean yield from 4,061 to 4,248 kg ha$^{-1}$; the paired differences became $-0.026$ ($-92$ kg ha$^{-1}$) against the rainfed October 15 schedule and $-0.126$ ($-656$ kg ha$^{-1}$) against the optimized fixed rule. The $k$NN rule rose to 1.081 (4,706 kg ha$^{-1}$, 26.3 mm). The result in 2024 did not change, because the 2024 context was complete: three of the five seeds still applied no irrigation in the drought season. Because the same rule would also mask windows of training seasons (2006, 2014 and 2016), a complete correction would require retraining; the analysis is reported as a sensitivity and not as the primary result.

%%============================================================
\subsection{Run integrity and scope of inference}
\label{subres_generalization}

No DSSAT run failed in the 9,030 development and 2,408 test grid simulations or in the PPO evaluations. No selected PPO action reached the bounds of the trigger range, and no evaluated depth fell below the 5 mm event threshold. Two late sowings selected by PPO in 2021 and 2022 (December 2 and 3) exceed by one and two days a conservative 150-day horizon ending on April 30; all other sowings remained inside it. The conclusions are restricted to Castro, to the synthetic profile and the generic cultivar, and to the irrigation-rule family described in Section~\ref{subdecision}.

% =====================================================================
\section{Discussion}
\label{discussao}

The main result is negative for the learned policy and positive for the comparator structure. Inside the training seasons PPO learned a policy that differed between seasons and exceeded the best fixed rule; on the validation seasons it also exceeded the fixed rule; on four later seasons it produced a lower objective and less yield than every comparator, and the optimized fixed rule captured two thirds of the gain available over the rainfed calendar. The reversal is explained by the information given to the policy and by the seasons used to select it, both of which can be checked with data available before the test block was opened.

%%============================================================
\subsection{Mechanisms: climate--soil--sowing--water interaction}
\label{subdisc_mechanisms}

In the simulated system, two processes determine the value of a management action. The sowing date places the reproductive period in a window of radiation, temperature and rainfall, and its effect is present in every season (a mean within-season range of 1,853 kg ha$^{-1}$). Irrigation has value only in the few seasons in which the 198 mm of available water of the profile is depleted during the reproductive period: three of fifteen development seasons and one of four test seasons. A fixed rule with a high trigger exploits this structure without knowing in advance which season will be dry, because the trigger reacts to the in-season water balance; in the test block it applied water only in 2023 (15 mm) and 2024 (75 mm). The contextual policies had to anticipate, from information available on September 1, whether the season would need water, and the correlations of the development seasons showed that the pre-season variables carried almost no information on the irrigation gain. This is consistent with the observation of \citet{kelly2024assessing} that optimized heuristic irrigation policies were difficult to outperform outside strongly water-limited conditions.

%%============================================================
\subsection{Value of context and information}
\label{subdisc_information}

The decomposition attributes the test gain to the choice of a good fixed rule and a negative contribution to context. Two features of the information explain this result. First, the context had low skill for the outcome that mattered most in the test block. The ONI terms, the only seasonal-scale information in the context, were not associated with sowing response or irrigation gain in the fifteen development seasons, which agrees with the view that the value of pre-season climate information must be measured for the decision at hand rather than assumed \citep{hansen2002climate}. Second, the context inherited the defects of the station record. Two test seasons had no valid temperature in the 30-day window, the interpolated values fell far outside the training range, and every seed responded with the latest sowing dates observed in the study. A neural policy extrapolates without signalling that its input is unfamiliar; the $k$NN rule, which can only reuse actions from observed seasons, was less affected in these two seasons. When the missing windows were treated as unknown, the late sowings disappeared and the gap to the rainfed calendar fell from 279 to 92 kg ha$^{-1}$.

The learning curves add a third element. With eleven training seasons the policy can associate each context with the best action of that season. The development run without regularization did so almost completely (training objective 1.187 against an oracle of 1.213) and lost on validation; the regularized runs generalized to validation, but the checkpoint was selected on four seasons in which no irrigation was needed. The validation objective therefore rewarded policies that did not irrigate, and the selected checkpoints, taken at 6,144--12,288 steps, applied no water in 2024 in three of five seeds. The composition of the validation block thus acted as an implicit prior on the irrigation behaviour of the policy.

%%============================================================
\subsection{When an adaptive policy is worth the complexity}
\label{subdisc_value}

The results suggest three conditions for a context-conditioned policy to add value over a fixed rule in a system of this type. The context must have measurable skill for the component of the decision that varies between seasons, here the occurrence of a reproductive-period drought; the training and validation seasons must include enough seasons of each hydrological type for the selection criterion to reward the right behaviour; and the pre-season inputs must pass a coverage check before they reach the policy, with a fallback when they do not. Where these conditions are not met, a fixed rule whose irrigation component reacts to the season, selected on all available development seasons, is the stronger comparator and should be reported in any evaluation of reinforcement learning for crop management \citep{balderas2025comparative, noa2026hybrid}. Options that address these conditions without using the test block include longer training records or weather generators to increase the number of distinct seasons, cross-validation over development seasons for checkpoint selection, and policies that restrict the action to the rules supported by the nearest observed contexts.

%%============================================================
\subsection{Resource allocation and risk}
\label{subdisc_resources}

The productivity--water relation of the test block had a steep initial segment: the fixed rule used 22.5 mm on average and raised mean yield by 564 kg ha$^{-1}$ over the rainfed calendar, almost entirely through 2024. Beyond about 40 mm, additional water did not raise mean yield. The low-yield tail was governed by 2022, a season in which no rule produced more than 3,139 kg ha$^{-1}$, and by the drought of 2024 for the rules that did not irrigate. PPO had the lowest P10 and CVaR$_{10}$ because its late sowings in 2022 added to the low yield of that season. These are properties of the simulated system with a generic cultivar and a synthetic soil, and they are not an operational recommendation of irrigation depth.

%%============================================================
\subsection{Scope and limitations}
\label{subdisc_limitations}

The scope of inference is Castro only. Algorithmic replication was tested with five seeds, but the agronomic sample remains four independent seasons in the test block, and one of them (2024) determined most of the difference between rules. The main limitations are the synthetic soil and generic cultivar without field calibration, the use of municipal yield only as contextual information, the quality of the station record in recent years, the proxy water balance of the irrigation rule, and the simplified one-step representation of management. The test block had been evaluated once in an earlier configuration with a different context and objective; the configuration reported here was fixed on the development seasons before its evaluation. Logistic constraints, irrigation costs beyond the linear penalty, nutrient limitations, cultivar choice, and machinery availability were not modelled.

% =====================================================================
\section{Conclusion}
\label{conclusao}

The study asked whether a context-conditioned sowing-and-irrigation policy adds value over a fixed rule in Castro. With raw DSSAT yield and information available on September 1, PPO learned season-specific policies that exceeded the optimized fixed rule on the eleven training seasons (objective 1.156 vs 1.130) and on the four validation seasons (1.226 vs 1.175). On the 2021--2024 test block the five PPO policies produced 4,061 kg ha$^{-1}$ on average, against 4,904 kg ha$^{-1}$ for the fixed rule optimized on all development seasons, a paired difference of $-843$ kg ha$^{-1}$ that was negative in every season. The optimized fixed rule captured 67\% of the gain of the retrospective oracle over the rainfed October 15 calendar, mainly by irrigating 75 mm in the 2024 drought.

Regarding the three directional expectations: the value of adaptation depended on water availability, because the only season with a large difference between rules was the one in which the soil water was depleted; the effect of soil storage appeared through the high trigger selected for the fixed rule, which irrigated only when the 198 mm profile was depleted; and the value of pre-season information was low, in line with its weak association with sowing response and irrigation gain in the development seasons. The analysis identified two causes of the out-of-sample loss that can be checked before a policy is deployed: pre-season inputs reconstructed from missing records, which moved two test contexts far outside the training range, and a validation block without dry seasons, which selected policies that did not irrigate. The contribution is a reproducible DSSAT--PPO workflow with a comparator structure that exposes these effects, and the evidence that, for this system and data, the design of the information set and of the validation seasons matters more than the choice of the learning algorithm. Because DSSAT was not calibrated with local field data, the results should not be interpreted as field-scale recommendations for sowing date or supplemental irrigation.
%===============================================================
\section*{Declaration of generative AI and AI-assisted technologies in the writing process}

The authors used ChatGPT only to support English-language editing and readability checks. The authors reviewed and verified the manuscript content, numerical results, cited literature, code outputs, and interpretations, and take full responsibility for the submitted version.

\section*{Declaration of competing interest}
The authors declare that they have no known financial or personal competing interests that could have influenced the work reported in this article.

\section*{Compliance with ethical standards}
This study did not involve human participants, animals, or sensitive personal data. Therefore, no ethical approval was required.

\section*{Data and code availability}

Code, data, figures, and accompanying materials for the reported computational workflow are available on Zenodo as software, version v1.0.0, at DOI \url{https://doi.org/10.5281/zenodo.22260538}. The repository is available at \url{https://github.com/GabyQueiroz/RLDSSAT_Soybean} and includes the executable computational environment, redistributable data, seeds, logs, scripts used to generate figures and tables, baseline configurations, and license metadata.

\section*{Funding}
This research was funded by national funds through FCT – Fundação para a Ciência e a Tecnologia, I.P., under project UIDB/00048/2020 (DOI: 10.54499/UIDB/00048/2020), by the National Council for Scientific and Technological Development (CNPq, Brazil) through Productivity in Research Scholarships (Grant Nos. 301644/2022-5 and 303909/2025-0), and by the Fundação Araucária through a postdoctoral fellowship (Grant No. 56/2026 FAUEPG).

\section*{CRediT authorship contribution statement}

Marcella Scoczynski Ribeiro Martins: Conceptualization, Investigation, Writing -- original draft, Writing -- review and editing. Gabrielly de Queiroz Pereira: Data curation, Formal analysis, Software, Visualization, Writing -- original draft, Writing -- review and editing. Hilkija Gaius Tosso: Methodology, Validation, Writing -- review and editing. Maria Salete Marcon Gomes Vaz: Supervision, Methodology, Writing -- review and editing. Jose Carlos Ferreira da Rocha: Supervision, Methodology, Writing -- review and editing. Viviane M. Gomes Pacheco: Methodology, Writing -- review and editing. Cloves Goncalves Rodrigues: Methodology, Writing -- review and editing. Antonio Paulo Coimbra: Supervision, Writing -- review and editing. Wesley Pacheco Calixto: Conceptualization, Methodology, Supervision, Project administration, Funding acquisition, Writing -- review and editing.

\section*{Supplementary material}

The supplementary material contains the source tables, diagnostic outputs, and additional analyses that support the results reported in the main text. The main article retains the system and decision-timeline figures, the DSSAT configuration and validation limits, the paired effects against the final comparators, the productivity--water analysis, the temporal inference scope for Castro, and the context--action relationship.

\begin{itemize}
\item \textbf{S1. PPO configuration.} Configuration file \texttt{experiment\_castro\_t0.yaml}, per-seed generated configurations, and the context scaler of each seed.
\item \textbf{S2. Soil profile.} Soil profile used by the executable DSSAT workflow, reported in \texttt{synthetic\_soil\_profile\_s2.csv}.
\item \textbf{S3. Weather, temporal split, and planting-window checks.} Weather completeness, imputation diagnostics, SIDRA pairing audit, temporal registry, ZARC audit, and planting-window checks.
\item \textbf{S4. Previous diagnostic configuration.} Statistical yield correction, its ablation, and the policy evaluated with realized-season context, retained to document the redesign.
\item \textbf{S5. Learning curves.} Validation curves, TensorBoard scalars, selected checkpoints and stopping steps of seeds 42--46, and the development runs of seed 42.
\item \textbf{S6. Candidate grid.} Simulated yield, irrigation and objective of the 602 grid rules in the development and test seasons.
\item \textbf{S7. Actions, paired effects, and productivity--water analysis.} Actions by season and seed, paired effects against all comparators, value decomposition, regret, and water-penalty sensitivity.
\item \textbf{S8. Context analysis.} Context variables by season, Spearman correlations with season outcomes, position of each context relative to the training seasons, and the post hoc coverage sensitivity.
\end{itemize}


\bibliographystyle{elsarticle-harv}
\bibliography{01_marcella_v3.bib}


% =====================================================================



\end{document}
